% Limitations (F18, F20, F24, F23, F22, F4, F19 traceable).
% Scale: 11M-124M (completed), 1.7B (in progress); Warstadt2023 BabyLM source.
% Step budget confound: 11M:20k vs 121M:10k both arms (F20); evaluated via Anderson (F24 Tarka scope).
% Judge bias: F23 per-domain decomposition (judge trained on childes-heavy strict-small).
% BLiMP resolution ~1.4pp at n=1000; single benchmark limitation.
% Deterministic allocator peak (F4); two-specialist restriction (F22).

\section{Limitations}
\label{sec:limitations}

The empirical program operates within the constraints of its pre-registered hypothesis gates and the available compute resources. All completed experiments (hypotheses H1–H4 and H6) span from 11~million to 121~million parameters, each trained on the 10~million-word BabyLM strict-small corpus~\cite{Warstadt2023}. Conversion of larger pretrained models (a 1.7~billion-parameter Qwen instance per F25) is underway but not complete; the findings therefore make claims about small-LM equilibrium dynamics grounded in experiments at these scales, not extrapolations to foundation models. The anytime training trajectory (F24, section~\ref{sec:h6}) includes unrolled supervision that achieves statistical parity with a 12-layer explicit baseline; that result is established at smoke (204-dimensional) and full 121~million-parameter scales, and does not yet apply to the larger-scale conversions where the solver-contraction challenge may re-emerge.

The step-budget confound between the 11~million and 121~million matched comparisons (F18 and F20) requires explicit treatment because it couples with the width scaling. The 11M baseline consumed 20,000 training steps; the 121M baseline consumed 10,000 steps. This budget difference was held constant across both arms (explicit and equilibrium architectures) within each scale study, so the direction of any quality gap cannot be reversed by budget alone—both models moved through identical step counts, and a quality reversal would require the equilibrium architecture to be simultaneously hurt more by width and helped more by fewer steps than the explicit model, an unlikely coincidence. However, the magnitude of the quality degradation (ratio 0.930 at 11M, ratio 0.787 at 121M) conflates the effects of width and budget. Isolating the width component from the step-count component would require a matched-step run at both scales or a matched-compute run where the explicit model receives proportionally more steps to balance wall-clock parity. The trend is therefore reported as a phenomenon (the quality gap widens with width) without a calibrated functional form, and the open problem is posed accordingly in section~\ref{sec:discussion}.

A second technical scoping applies to the evaluation path and the nature of the parity claim. The anytime training arms (F24, B1) supervise the model on plain forward-pass iterates $z_4, z_8, z_{12}$ of the post-LN fixed-point map. Evaluation of these checkpoints applies Anderson-accelerated fixed-point solving at the same iteration budgets (4, 8, 12)—a different computational procedure. The parity result (BLiMP ratio 0.991 at 121M) and the graceful degradation across budgets both attach to the Anderson-evaluated path, not to equivalence under plain iteration. This distinction proved critical in follow-up generative experiments (F24 addendum): one checkpoint generating coherent English through unrolled computation produced degenerate punctuation loops through the Anderson solver, a failure mode isolated to Anderson and resolved by matching the checkpoint's decode mode to its training procedure. The fix ensures the parity claim is evaluated consistently, but it narrows the scope: the finding is that anytime supervision recovers the explicit baseline's quality under Anderson evaluation, not that the two algorithms are functionally identical.

The closed-loop generation verdict (F23, formally MISSED) is specific to the evaluation judge and does not generalise to arbitrary models or judges. The held-out test used a frozen 121~million-parameter explicit language model trained on BabyLM strict-small to score generated continuations by negative log-likelihood. This judge exhibits a systematic style bias: continuations generated in the childes (child-directed speech) register score approximately 2~nats/token better than equivalently capable text in the simple-wiki register, independent of which system produced the output. The per-domain decomposition (F23, Tarka-directed rerun) quantifies this confound: the best specialist on wiki-only prompts scores 3.27–3.78 nats/token when generating wiki-style and 5.10–5.49 when forced to generate childes-style, a gap exceeding the 1.5–2~nat effect size of the closed-loop generation loss. The auction's true advantage or disadvantage under a style-neutral judge remains unknown; no unbiased judge at this model scale exists in-repository, and querying an external LLM would introduce a different source of bias (model choice, scale, training data). The claim that the auction loses under the BabyLM-distribution judge is precise and data-supported, but it is narrower than the pre-registered hypothesis, which made no judge specification.

BLiMP serves as the sole linguistic benchmark applied to all models. The evaluation is restricted to 1,000 scored pairs (a subset of the full corpus), and the statistical resolution of this benchmark—derived from a binomial model—is approximately 1.4 percentage points at the 95\% confidence level. The pre-registered H1 threshold of 95\% parity sits within this margin of resolution, and the observed ratio (0.930, 95\% CI [0.898, 0.949]) crosses the threshold only in its lower confidence bound. A complementary benchmark, such as GLUE or SUPERGLUE subsets or manual evaluation of generative quality, would strengthen the evaluation; the present results measure preference on one structured grammaticality task. The use of BLiMP alone does not invalidate the findings, but it represents a single-task evaluation whose findings may not transfer to other dimensions of language understanding.

The peak activation memory reported for both small and large scales (F4, F20) reflects the architecture's deterministic allocator peak rather than a stochastic measurement across independent runs. For the 121~million equilibrium models, peak memory is 6.29~gigabytes, bit-identical across seeds 42, 43, and 44. This exact reproducibility indicates that the peak is determined by the static computational graph (the number of solver iterations, activation storage pattern) rather than varying with initialization or random sampling during training. The memory advantage is therefore an architecture-level property and no confidence interval is computed; the claim is that the equilibrium depth-memory relationship holds for the specified graph structure, not that per-run variability is negligible.

Finally, the token-auction studies (F22, F23) operate on exactly two 30~million-parameter domain specialists (one trained on childes, one on simple-wiki) aggregated via second-price mechanism. The mechanism's truthfulness (F6, empirical regret exactly 0.0) and the per-token selection advantage at scoring time (23\% ppl improvement, F22) are validated for this two-specialist configuration. Whether the benefits persist with larger expert sets (three, five, or more specialists), with finer-grained domain decomposition, or with learned bid functions rather than confidence-only bids remains untested. The confidence-bid restriction is theoretically sound for auction mechanism design but limits the model's adaptability: specialists bid only their own max-probability confidence, not the prefix's measured difficulty or domain, so the mechanism cannot bias selection toward the specialist most suited to the current context in a principled way. Addressing this limitation requires a multi-input learned bidding function and revalidating truthfulness under that new mechanism—a follow-up program.
